% =============================================================================
% 035 / 068
% Status: raw
% =============================================================================
% Notes:
%
% Source question:
% 閣下整理出吳語播音體欲以為吳語書面語，那閣下是如何看待那些並沒有那麼多時態等的吳語方言的呢？
% =============================================================================

\chapter[閣下整理出吳語播音體欲以為吳語書面語，那閣下是如何看待那些並沒有那麼多時態等的吳語方言的呢？]{閣下整理出吳語播音體欲以為吳語書面語，那閣下是如何看待那些並沒有那麼多時態等的吳語方言的呢？}
\qaanswer{
幾乎所有的吳語亞方言都或多或少有絕對時態標記(時態)和相對時態標記(事態)。但吳語區的府城被官話和普通話滲透、漂變過深，很多方面已經藍青化和普化，自然就看不到時態的痕跡了。所以問題在於，吳語構建按官話，普通話的語法來合適還是回歸原生吳語本源合適？走前者的路線短期無需支付大的成本但最終會使吳語徹底消亡，走後者的路線成本大且艱辛，但能讓吳語作為一種獨立語言延續下去，最重要的是能將吳語語言思維這種獨特的思維方式延續下去。
}

